
\documentclass{article} % Clase base sin tamaño

% Página horizontal
\usepackage[a4paper,landscape,margin=2cm]{geometry}

% Arial y tamaño exacto
\usepackage{fontspec}
\setmainfont{Arial}
\usepackage{xcolor}
% Para tamaño grande explícito
\usepackage{setspace}
\usepackage{graphicx}

\usepackage{amsmath}
\usepackage[thmmarks, amsmath, thref]{ntheorem}
\usepackage{amsfonts}
\usepackage{amssymb}
\usepackage{polyglossia}       % Paquete multilingüe moderno
\setmainlanguage{spanish}   

\newcommand{\bigtable}[1]{{\fontsize{30pt}{36pt}\selectfont\renewcommand{\arraystretch}{2}#1}}


\begin{document}

% TODO lo que esté entre estas llaves tendrá tamaño 30pt real
{\fontsize{30pt}{36pt}\selectfont

\begin{center}
\
\textbf{\\CERTAMEN 1}
\\
\end{center}

\begin{center}

\textbf{Econometría ICS-294}
\\
\vspace{0.2cm}
I Semestre 2025\\
Profesores: Jocelyn Tapia Stefanoni, Franco Antonucci Fernández y Ricardo Pruneda Urz\'ua\\
%Ayudante: 
\end{center}

\textbf{Instrucciones:} 

\begin{itemize}

    \item Usted dispone de 90 minutos  para intentar los 100 puntos del certamen.
    \item Cuide la presentación y redacción de sus preguntas.
    \item No puede utilizar celular, computador o tablet, pero sí una calculadora TI.
     
\end{itemize}

{\bf RdA3.2  Preguntas cortas (20 puntos)}\\

{\bf Elija solo 2 preguntas y responda. }
\begin{enumerate}
    \item (10 puntos) Dos economistas, Milana y Maxi, están analizando datos de una muestra de personas de todos los países de América Latina para estudiar cómo influye la \textit{cantidad de aplicaciones abiertas en el celular} sobre la \textit{duración de la batería}. 
    Debido a posibles errores de medición en el proceso de recolección de datos, surge un debate entre ambas. Milana sostiene que, desde el punto de vista econométrico y en términos de obtener un mejor estimador, sería menos problemático que el error de medición estuviera en la variable \textit{duración de la batería}. Por su parte, Maxi argumenta que, por el contrario, sería menos grave que el error estuviera en la variable \textit{cantidad de aplicaciones abiertas en el celular}. ¿Quién tiene la razón? ¿Por qué? Puede apoyarse con herramientas matemáticas.

 
    %intervalos de confianza
    \item (10 puntos) Describa el esquema de muestreo utilizado por Miguel y Kremer (2004) en su estudio {\it Worms: Identifying Impacts on Education and Health in the Presence of Treatment Externalities}.
    ¿Qué ventajas tiene este diseño para la identificación causal de los efectos del tratamiento?

\item  (10 puntos) Josefina se encuentra trabajando para una importante constructora nacional en el \'area de finanzas, el jefe de
proveedores le comenta que la media de cuentas por cobrar es de 220.000.000 de pesos. Josefina tiene una  muestra aleatoria de 500 facturas con media igual a 195.000.000 de pesos y una desviaci\'on
est\'andar de 1.000.000. Con un 90\% de confianza ¿es correcta la afirmacion del jefe?


%     $$\color{red}{\bar{x} \pm z_{\frac{\alpha}{2}}\frac{\sigma}{\sqrt{n}}}$$

% {\color{red}$[194.926.209,8 ; 195.073.790,2]$}


%   {\color{red}  El valor que comenta el jefe es 220.000.000 el cual no está en el intervalo de confianza calculado en a), así que  90\% de confianza la afirmación es incorrecta.}

\item (10 puntos) La  suma de cuadrados de los residuos es:
    \[
    SSR = \sum_{i=1}^{n} \hat{u}_i^2 = \hat{u}^T \hat{u}
    \]
    
Demuestre que esta puede reescribirse como:
    \[
    SSR = y^T y - \hat{y}^T y
    \]
    donde \( \hat{y} = X \hat{\beta} \) representa los valores ajustados del modelo, y \( \hat{u} = y - \hat{y} \) es el vector de residuos.

\end{enumerate}

\vspace{1.5cm}

{\bf RdA3.2 Ejercicios (80 puntos)}\\

{\bf Elija 4 de las siguientes preguntas y responda.} 
\begin{enumerate}
  \item La gerencia de Recursos Humanos de una reconocida empresa de servicios está interesada en mejorar la calidad de la atención mediante la introducción de una nueva herramienta de trabajo. Para ello, ha invitado a los trabajadores con más de 2 años en la empresa a inscribirse voluntariamente en una capacitación. De un total de 15 empleados, 5 decidieron capacitarse. Un año después, el salario promedio del grupo capacitado asciende a US\$1.850, mientras que el del grupo que no se capacitó es de US\$1.600.
\begin{enumerate}
      \item  (5 puntos) ¿Podemos concluir que existe una relación \textit{causal} entre participar en la capacitación y el aumento del salario? Justifique.

    \item  (7 puntos) Halle el valor del coeficiente $\hat{\beta}$ en un modelo de regresión lineal simple que explique el salario en función de la variable dicotómica ``capacitación'' (1 = capacitado, 0 = no capacitado).
    \item  (8 puntos) Se sabe que las \textit{desviaciones estándar} del salario para el grupo capacitado y para el no capacitado  son  US\$250  y  US\$200 respectivamente. Con un nivel de confianza del 95\%, ¿es estadísticamente significativa la diferencia de medias entre ambos grupos?
\end{enumerate}
    
%     \item   Dado el siguiente modelo de regresión poblacional, $Y_i = \beta_0 + \beta_1 X_i + u_i$.
%     \begin{enumerate}
%         \item (10) Si usted cuenta con una muestra de tamaño n. Pruebe que, los estimadores de los parámetros $\beta_0$ y $\beta_1$ calculados por el m\'etodo de  MCO son:\\
%         $\hat{\beta_1}=\frac{\sum_{i=1}^{n}(x_i - \bar{x})(y_i-\bar{y})}{\sum_{i=1}^{n}(x_i - \bar{x})^2}$ y  $\hat{\beta_0}=\bar{Y}-\hat{\beta}_1\bar{X}$.
%        \item (10 puntos)	¿Qu\'e supuestos se deben satisfacer para que los estimadores  sean insesgados? Asumiendo que se satisfacen los supuestos muestre que son insesgados. \end{enumerate}
%      \item (20 puntos) Sea un modelo de regresión lineal simple en el cual, en la práctica, se observan $
% Y^* = Y + u, \quad \text{donde } u \sim \mathcal{N}(0,\,1),
% $ y $
% X^* = X + c$, \quad \text{donde } $c \neq 0$ \text{ es una constante}. Si el modelo verdadero (sin errores de medición) es $
% Y = \beta_0 + \beta_1 \, X + \varepsilon$, compare los estimadores obtenidos  usando \((Y^*,\,X^*)\) versus  \((Y,\,X)\).





\item (20 puntos) Considere una muestra de tamaño n, pero en la cual las variables observadas presentan errores de medición:

\[
Y_i^* = Y_i + u_i, \quad \text{donde } \quad u_i \sim \mathcal{N}(0,\,1), \quad \text{y}
\]
\[
X_i^* = X_i + c, \quad \text{donde } \quad c \neq 0 \quad \text{ es} \quad  \text{constante}
\]

El modelo verdadero (sin errores de medición) es:
\[
Y_i = \beta_0 + \beta_1 X_i + \varepsilon_i
\]

Compare los estimadores de Mínimos Cuadrados Ordinarios (MCO) obtenidos utilizando los datos observados \((Y_i^*, X_i^*)\) versus aquellos que se obtendrían si se observaran las variables verdaderas \((Y_i, X_i)\), en una muestra de tamaño n. Analice el efecto que tienen los errores de medición en la estimación de los parámetros del modelo.




\item Considere el siguiente modelo de regresión poblacional:

\[
Y_i = \beta_0 + \beta_1 X_i + \varepsilon_i
\]

\begin{enumerate}
    \item[(a)] (10 puntos) Suponiendo que se dispone de una muestra aleatoria de tamaño \( n \), demuestre que los estimadores  MCO para los parámetros \( \beta_0 \) y \( \beta_1 \) están dados por:
    \[
    \hat{\beta}_1 = \frac{\sum_{i=1}^{n}(X_i - \bar{X})(Y_i - \bar{Y})}{\sum_{i=1}^{n}(X_i - \bar{X})^2}, \qquad
    \hat{\beta}_0 = \bar{Y} - \hat{\beta}_1 \bar{X}
    \]

    \item[(b)] (10 puntos) Indique y explique los supuestos necesarios para que los estimadores MCO sean insesgados. Luego, bajo dichos supuestos, pruebe que los estimadores de \( \beta_0 \) y \( \beta_1 \) son efectivamente insesgados.
\end{enumerate}

\vspace{0.5cm}




    \item Considere el siguiente modelo de regresión lineal con los supuestos estándares:

    \[
    Y_i = \beta_0 + \beta_1 X_{1i} + \beta_2 X_{2i} + u_i
    \]
    
    Se dispone de la siguiente información:
    
    \[
    X'X = \begin{pmatrix}
    33 & 0 & 0 \\
    0 & 40 & 20 \\
    0 & 20 & 60
    \end{pmatrix}, \quad
    X'Y = \begin{pmatrix}
    132 \\
    24 \\
    92
    \end{pmatrix}
    \]
    
    \[
    \sum_{i=1}^{N} Y_i^2 = 678, \quad \sum_{i=1}^{N} (Y_i - \bar{Y})^2 = 150
    \]

%\begin{comment}
    \begin{enumerate}
        \item (3 puntos) ¿Cuál es el tamaño de la muestra?\\
    
        \item (9 puntos) Calcule los estimadores de MCO del modelo anterior.\\
        \textbf{Hint:} \( Z = \begin{pmatrix} A &{\bf 0} \\ {\bf 0} & B \end{pmatrix} \) con \( A \) y \( B \) matrices cuadradas, entonces:\[Z^{-1} = \begin{pmatrix} A^{-1} & {\bf 0} \\{\bf 0} & B^{-1} \end{pmatrix}\]

        \item (4 puntos) Indique la ecuación del modelo predicho e interprete.

        \item (4 puntos) ¿Qué debería suceder con $R^2$ si agregamos más variables independientes a la medición?
        
        % \item ¿Qué valor alcanza el $R^2$?\footnote{Hint: Recuerde que $SSE = \sum_{i=1}^N (\hat{Y_i} - \bar{Y})^2$ que de manera matricial queda como $SSE=\hat{Y}^T Y - N \bar{Y^2}$}
        %regresion lineal multiple
    \end{enumerate}

\item  Supongamos que eres un investigador interesado en estudiar la relación entre la cantidad de horas de estudio (variable independiente X) y el puntaje obtenido en un examen (variable dependiente Y).
Tienes los siguientes datos de 6 estudiantes:

\bigtable{
\begin{tabular}{|c|c|}
\hline
\textbf{Horas de estudio (X)} & \textbf{Puntaje en el examen (Y)} \\
\hline
2 & 50 \\
3 & 60 \\
4 & 65 \\
5 & 70 \\
6 & 75 \\
7 & 76 \\
\hline
\end{tabular}
}

\begin{enumerate}
    \item  (5 puntos) Use el método de los mínimos cuadrados para ajustar una línea recta $\hat{Y}=\hat{\beta_0}+\hat{\beta}_1X$.
   \item (5 puntos) ¿Cu\'al es el puntaje esperado en el examen para un estudiante que estudia 8 horas?
\item (5 puntos) Si un estudiante aumenta sus horas de estudios en 2 horas ¿cu\'anto se espera que aumente su puntaje?
\item (5 puntos) Calcule el $R^2$ e interprete.
\end{enumerate} 

\end{enumerate}

   \newpage 
{\bf Formulario}
\begin{itemize}

\item El intervalo de confianza de $100(1-\alpha)\%$ para la media poblacional $\mu$ cuando $\sigma$ es \emph{\textbf{conocido}}:
$\bf {\bar{x} \pm z_{\frac{\alpha}{2}}\frac{\sigma}{\sqrt{n}}}$

\item El intervalo de confianza de $100(1-\alpha)\%$ para la diferencia de medias
$\mu_x-\mu_y$ es:
${\bf(\bar{X}-\bar{Y}) \pm z_{\frac{\alpha}{2}}\sqrt{\frac{\sigma_x^2}{n_x}+\frac{\sigma_y^2}{n_y}}}$

\item  Estimadores MCO.
{\bf
$$\hat{\beta}_0 = \bar{y} - \hat{\beta}_1 \bar{x} \quad \quad   \hat{\beta}_1 = \frac{\sum_{i=1}^{n}(x_i - \bar{x})(y_i-\bar{y})}{\sum_{i=1}^{n}(x_i - \bar{x})^2}$$}
    
\item 
	{\bf	\[ SST = \sum_{i=1}^{n} (y_i - \bar{y})^2 \quad \quad  SSE = \sum_{i=1}^{n} (\hat{y}_i - \bar{y})^2 \quad \quad SSR = \sum_{i=1}^{n} \hat{u}_i^2  \]}
		
		
    \item Vector de estimadores:    
{\bf   $\hat{\beta}=(X^TX)^{-1}X^Ty$}

\item {\bf\[
A = \begin{bmatrix}
a & b \\
c & d
\end{bmatrix}
\qquad
A^{-1} = \frac{1}{det(A)} \begin{bmatrix}
d & -b \\
-c & a
\end{bmatrix}
\]}
\end{itemize}

}

\begin{figure}
    \centering
    \includegraphics[width=1.2\linewidth]{certamenes/normal_re.png}
\end{figure}
\end{document}
